\begin{table}[htbp]
\normalsize
\renewcommand{\arraystretch}{1.25}
\setlength{\tabcolsep}{6pt}
\centering
\caption{Individual characteristics and migration intention: linear probability models}
\label{tab:individual_characteristics_lpm_stata}
\setlength{\tabcolsep}{4pt}
\captionsetup{justification=centering}
\begin{tabular}{lccc}
\hline
 & (1) & (2) & (3) \\
\hline
Age & -0.006*** & -0.005*** & -0.005*** \\
 & (0.000) & (0.000) & (0.000) \\
Male & 0.051*** & 0.048*** & 0.053*** \\
 & (0.012) & (0.011) & (0.011) \\
Unemployed & 0.030 & 0.040 & 0.035 \\
 & (0.027) & (0.024) & (0.023) \\
Partnered & -0.051*** & -0.044*** & -0.056*** \\
 & (0.011) & (0.009) & (0.010) \\
High School or more & 0.026** & 0.021** & 0.032*** \\
 & (0.011) & (0.010) & (0.010) \\
Left-right ideology & 0.002 & 0.002 &  \\
 & (0.002) & (0.002) &  \\
Very interested in politics & 0.018 & 0.025* & 0.016 \\
 & (0.013) & (0.013) & (0.013) \\
Blank/null vote & 0.036** &  & 0.036*** \\
 & (0.015) &  & (0.013) \\
\hline
Observations &        4,966 &        6,155 &        5,641 \\
Year FE & Yes & Yes & Yes \\
Municipality FE & Yes & Yes & Yes \\
Survey weights & Yes & Yes & Yes \\
Blank/null vote included & Yes & No & Yes \\
Ideology included & Yes & Yes & No \\
\hline
\end{tabular}
\par\vspace{1.5mm}
\begin{minipage}{0.98\linewidth}
\footnotesize\setlength{\parindent}{0pt}\setlength{\parskip}{1pt}
Notes: The dependent variable is migration intention. Entries are LPM coefficients, with standard errors clustered at the municipality level and reported in parentheses. Each model uses exactly the estimation sample of its corresponding Logit specification and the same survey weights and fixed effects. The unweighted percentages of fitted values outside [0,1] are  9.69\%,  9.88\%, and 10.21\% in columns (1), (2), and (3), respectively. * \(p<0.10\), ** \(p<0.05\), *** \(p<0.01\).
\end{minipage}
\end{table}
